\BourbakiExercisesSection

\begin{enumerate}
\def\labelenumi{\arabic{enumi}.}
\item
  设 E, F 为两个 A-模，且 \(u: \mathbf{E} \to \mathbf{F}\) 为线性映射。设 M 为 E 的子模，N 为 F 的子模， \(i: \mathbf{M} \to \mathbf{E}\) 为典范单射， \(q: \mathbf{F} \to \mathbf{F} / \mathbf{N}\) 为典范满射。证明 \(u(\mathbf{M}) = \operatorname{Im}(u \circ i)\) 和 \(u^{-1}(\mathbf{N}) = \operatorname{Ker}(q \circ u)\) .
\item
  设 E 为左 A-模；对于 A 的每个左理想 a，令 aE 表示所有 ax 的和，其中 \(x \in E\) ，因此它是 E 的子模。
\end{enumerate}

\begin{enumerate}
\def\labelenumi{(\alph{enumi})}
\item
  证明对于每个 A-模族 \((\mathbf{E}_{\lambda})_{\lambda \in L}\) ，a. \(\bigoplus_{\lambda \in L} \mathbf{E}_{\lambda}\) 典范地等同于 \(\bigoplus_{\lambda \in L} a E_{\lambda}\) .
\item
  若 \(\mathfrak{a}\) 是 \(\mathbf{A}\) 的有限生成左理想，证明对每个族 \((\mathbf{E}_{\lambda})_{\lambda \in L}\) （由 \(\mathbf{A}\) -模组成）， \(\mathfrak{a}\) . \(\prod_{\lambda \in L} \mathbf{E}_{\lambda}\) 典范地等同于 \(\prod_{\lambda \in L} \mathfrak{a}\mathbf{E}_{\lambda}\) 。给出一个例子，其中 \(\mathfrak{a}\) 不是有限生成的且上述性质不成立（取 \(\mathbf{A}\) 交换且所有 \(\mathbf{E}_{\lambda}\) 都等于 \(\mathbf{A}\) ）.
\end{enumerate}

*(c) 设 K 是一个交换域，A 是 K 上两个未定元的多项式环 K{[}X,Y{]} ，m = AX 且 n = AY。若 E 是 \(\mathbf{A}^2\) 关于 \(\mathbf{A}^2\) 中由 \(\mathbf{Xe}_1 - \mathbf{Ye}_2\) 生成的子模的商 A-模 (其中 \((e_1, e_2)\) 是 \(\mathbf{A}^2\) 的典范基 )，证明 \((m \cap n)E \neq (mE) \cap (nE)\) .*

\begin{enumerate}
\def\labelenumi{\arabic{enumi}.}
\setcounter{enumi}{2}
\tightlist
\item
  设 E 是一个 A-模。
\end{enumerate}

\begin{enumerate}
\def\labelenumi{(\alph{enumi})}
\tightlist
\item
  设 \((\mathbf{L}_{\mu})_{\mu \in \mathbf{M}}\) 是一族右有向有序集，且对所有 \(\mu \in \mathbf{M}\) ，设 \((\mathrm{F}_{\lambda \mu})_{\lambda \in L_{\mu}}\) 是 \(\mathbf{E}\) 的一族递增的子模。证明
\end{enumerate}

\[
\bigcap_ {\mu \in \mathbf {M}} \left(\sum_ {\lambda \in L _ {\mu}} F _ {\lambda \mu}\right) = \sum_ {\rho \in \prod_ {\mu} L _ {\mu}} \left(\bigcap_ {\mu \in \mathbf {M}} F _ {\rho (\mu), \mu}\right).
\]

\begin{enumerate}
\def\labelenumi{(\alph{enumi})}
\setcounter{enumi}{1}
\tightlist
\item
  设 A 为一个域，E 为 A-模 \(A^{N}\) ；设 G 为 E 的子模 \(\mathbf{A}^{(\mathbf{N})}\) ；另一方面，对 N 的每个有限子集 H，设 \(F_{H}\) 为子模 \(\prod_{n\in \mathbf{N}}\mathbf{M}_n\) （ \(\mathbf{E}\) 的），其中 \(\mathbf{M}_n = \mathbf{A}\) （对于 \(n\notin \mathrm{H}\) ）且 \(\mathbf{M}_n = 0\) （对于 \(n\in \mathrm{H}\) ）。证明族 \((\mathrm{F_H})\) （其中 \(\mathbf{H}\) 遍历集 \(\Phi\) ，由 \(\mathbf{N}\) 的有限子集组成）是右有向的，并且
\end{enumerate}

\[
\left(\bigcap_ {\mathrm{H} \in \Phi} \mathrm{F} _ {\mathrm{H}}\right) + \mathrm{G} \neq \bigcap_ {\mathrm{H} \in \Phi} (\mathrm{F} _ {\mathrm{H}} + \mathrm{G}).
\]

\begin{enumerate}
\def\labelenumi{\arabic{enumi}.}
\setcounter{enumi}{3}
\tightlist
\item
  设 \(\mathbf{E}, \mathbf{F}_1, \mathbf{F}_2\) 为三个 A-模， \(f_1: \mathbf{F}_1 \to \mathbf{E}\) , \(f_2: \mathbf{F}_2 \to \mathbf{E}\) 为两个 A-线性映射。 \(\mathbf{F}_1\) 和 \(\mathbf{F}_2\) 相对于 \(f_1\) 和 \(f_2\) 的纤维积是积 \(\mathbf{F}_1 \times \mathbf{F}_2\) 中由有序对 \((x_1, x_2)\) （满足 \(f_1(x_1) = f_2(x_2)\) ）组成的子模；记作 \(\mathbf{F}_1 \times_{\mathbf{E}} \mathbf{F}_2\) .
\end{enumerate}

\begin{enumerate}
\def\labelenumi{(\alph{enumi})}
\item
  对于每一对 A-线性映射 \(u_1 \colon G \to F_1, u_2 \colon G \to F_2\) （满足 \(f_1 \circ u_1 = f_2 \circ u_2\) ），存在且仅存在一个 A-线性映射 \(u \colon G \to F_1 \times_E F_2\) ，使得 \(p_1 \circ u = u_1, p_2 \circ u = u_2\) ，其中 \(p_1\) 和 \(p_2\) 表示 \(F_1 \times_E F_2\) 上投影 \(pr_1\) 和 \(pr_2\) 的限制.
\item
  设 \(\mathbf{E}'\) , \(\mathbf{F}_1'\) , \(\mathbf{F}_2'\) 是三个A-模， \(f_1': \mathbf{F}_1' \to \mathbf{E}', f_2': \mathbf{F}_2' \to \mathbf{E}'\) 为两个 A-线性映射，而 \(\mathbf{F}_1' \times_{\mathbf{E}'} \mathbf{F}_2'\) 为相对于这些映射的纤维积。对于每一组 A-线性映射 \(v_1: \mathbf{F}_1 \to \mathbf{F}_1'\) , \(v_2: \mathbf{F}_2 \to \mathbf{F}_2'\) , \(w: \mathbf{E} \to \mathbf{E}'\) ，若 \(f_1' \circ v_1 = w \circ f_1, f_2' \circ v_2 = w \circ f_2\) ，设 \(v\) 是 \(\mathbf{F}_1 \times_{\mathbf{E}} \mathbf{F}_2\) 到 \(\mathbf{F}_1' \times_{\mathbf{E}'} \mathbf{F}_2'\) 的A-线性映射，对应于两个线性映射 \(v_1 \circ p_1\) 和 \(v_2 \circ p_2\) 。证明如果 \(v_1\) 和 \(v_2\) 是单射，则 \(v\) 也是单射；给出一个例子，其中 \(v_1\) 和 \(v_2\) 是满射但 \(v\) 不是满射（取 \(\mathbf{E}' = \{0\}\) ）。如果 \(w\) 是单射且 \(v_1\) 和 \(v_2\) 是满射，证明 \(v\) 是满射。
\end{enumerate}

\begin{enumerate}
\def\labelenumi{\arabic{enumi}.}
\setcounter{enumi}{4}
\tightlist
\item
  设 \(\mathbf{E}, \mathbf{F}_1, \mathbf{F}_2\) 为三个 A-模， \(f_1: \mathbf{E} \to \mathbf{F}_1, f_2: \mathbf{E} \to \mathbf{F}_2\) 为两个线性映射。 \(\mathbf{F}_1\) 和 \(\mathbf{F}_2\) 相对于 \(f_1\) 和 \(f_2\) 的融合和是 \(\mathbf{F}_1 \times \mathbf{F}_2\) 关于 \(\mathbf{E}\) 在映射 \(z \mapsto (f_1(z), -f_2(z))\) 下的像所成的子模的商模；记作 \(\mathbf{F}_1 \oplus_{\mathbf{E}} \mathbf{F}_2\) .
\end{enumerate}

\begin{enumerate}
\def\labelenumi{(\alph{enumi})}
\tightlist
\item
  对于每一对 A-线性映射 \(u_{1}: \mathbf{F}_{1} \to \mathbf{G}, u_{2}: \mathbf{F}_{2} \to \mathbf{G}\) ，若 \(u_{1} \circ f_{1} = u_{2} \circ f_{2}\) ，证明存在且仅存在一个 A-线性映射 \(u: \mathbf{F}_{1} \oplus_{\mathbf{E}} \mathbf{F}_{2} \to \mathbf{G}\) ，使得 \(u \circ j_{1} = u_{1}, u \circ j_{2} = u_{2}\) ，其中 \(j_{1}\) 和 \(j_{2}\) 分别表示典范映射
\end{enumerate}

\[
\mathbf {F} _ {1} \times \mathbf {F} _ {2} \rightarrow \mathbf {F} _ {1} \oplus_ {\mathrm{E}} \mathbf {F} _ {2}
\]

与典范单射 \(\mathbf{F}_1\to \mathbf{F}_1\times \mathbf{F}_2,\mathbf{F}_2\rightarrow \mathbf{F}_1\times \mathbf{F}_2.\)

\begin{enumerate}
\def\labelenumi{(\alph{enumi})}
\setcounter{enumi}{1}
\tightlist
\item
  设 \(\mathbf{E}'\) , \(\mathbf{F}_1'\) , \(\mathbf{F}_2'\) 是三个A-模， \(f_1': \mathbf{E}' \to \mathbf{F}_1'\) , \(f_2': \mathbf{E}' \to \mathbf{F}_2'\) 为两个 A-线性映射，而 \(\mathbf{F}_1' \oplus_{\mathbf{E}'} \mathbf{F}_2'\) 为相对于这些映射的融合和。对于每一组 A-线性映射 \(v_1: \mathbf{F}_1' \to \mathbf{F}_1\) , \(v_2: \mathbf{F}_2' \to \mathbf{F}_2\) , \(w: \mathbf{E}' \to \mathbf{E}\) ，若 \(v_1 \circ f_1' = f_1 \circ w\) , \(v_2 \circ f_2' = f_2 \circ w\) ，设 \(v\) 是 \(\mathbf{F}_1' \oplus_{\mathbf{E}'} \mathbf{F}_2'\) 到 \(\mathbf{F}_1 \oplus_{\mathbf{E}} \mathbf{F}_2\) 的A-线性映射，对应于两个线性映射 \(j_1 \circ v_1\) 和 \(j_2 \circ v_2\) 。证明如果 \(v_1\) 和 \(v_2\) 是满射，那么 \(v\) 也是满射；给出一个例子，其中 \(v_1\) 和 \(v_2\) 是单射但 \(v\) 不是单射。证明如果 \(w\) 是满射且 \(v_1\) 和 \(v_2\) 是单射，则 \(v\) 是单射。
\end{enumerate}

\begin{enumerate}
\def\labelenumi{\arabic{enumi}.}
\setcounter{enumi}{5}
\tightlist
\item
  给定一个 A-模 E，设 \(\gamma(E)\) 表示 E 的生成系的基数中的最小者。
\end{enumerate}

\begin{enumerate}
\def\labelenumi{(\alph{enumi})}
\tightlist
\item
  若 \(0 \to \mathbf{E} \to \mathbf{F} \to \mathbf{G} \to 0\) 是 A-模的正合序列，则
\end{enumerate}

\[
\gamma (\mathbf {G}) \leqslant \gamma (\mathbf {F}) \leqslant \gamma (\mathbf {E}) + \gamma (\mathbf {G}).
\]

给出一个乘积 \(\mathbf{F} = \mathbf{E} \times \mathbf{G}\) 的例子，使得 \(\gamma(\mathbf{E}) = \gamma(\mathbf{F}) = \gamma(\mathbf{G}) = 1\) （取 \(\mathbf{E}\) 和 \(\mathbf{G}\) 为 \(\mathbf{Z}\) 的商模 ）.

\begin{enumerate}
\def\labelenumi{(\alph{enumi})}
\setcounter{enumi}{1}
\item
  若 \(\mathbf{E}\) 是子模族 \((\mathrm{F}_{\lambda})_{\lambda \in L}\) 的和，则 \(\gamma (\mathbf{E})\leqslant \sum_{\lambda \in L}\gamma (\mathbf{F}_{\lambda}).\)
\item
  若 \(\mathbf{M}\) 和 \(\mathbf{N}\) 是模 \(\mathbf{E}\) 的子模，证明
\end{enumerate}

\[
\sup (\gamma (\mathbf {M}), \gamma (\mathbf {N})) \leqslant \gamma (\mathbf {M} \cap \mathbf {N}) + \gamma (\mathbf {M} + \mathbf {N}).
\]

\begin{enumerate}
\def\labelenumi{\arabic{enumi}.}
\setcounter{enumi}{6}
\tightlist
\item
  \begin{enumerate}
  \def\labelenumii{(\alph{enumii})}
  \tightlist
  \item
    设 A 为一个环，E 为一个 (A, A)-双模；在积 \(B = A \times E\) 上，通过令下式成立来定义环结构
  \end{enumerate}
\end{enumerate}

\[
(a, x) (a ^ {\prime}, x ^ {\prime}) = (a a ^ {\prime}, a x ^ {\prime} + x a ^ {\prime});
\]

E（典范地等同于 \(\{0\} \times \mathbf{E}\) ）于是是B的一个双边理想，使得 \(\mathbf{E}^2 = \{0\}\) .

\begin{enumerate}
\def\labelenumi{(\alph{enumi})}
\setcounter{enumi}{1}
\tightlist
\item
  通过适当选取 A 和 E，证明 E 可以是一个左 B-模，其 \(\gamma(\mathbf{E})\) （习题6）为任意基数，尽管 \(\gamma(\mathbf{B}_{s}) = 1\) .
\end{enumerate}

*8. 设 A 是一个交换环， \(\mathbf{B} = \mathbf{A}[\mathbf{X}_n]_{n\geqslant 1}\) 是 A 上含无穷多个未定元的多项式环，C 是 B 关于由多项式 \((\mathbf{X}_1 - \mathbf{X}_2)\mathbf{X}_i\) （对于 \(i\geqslant 3\) ）生成的理想的商环。若 \(\xi_1\) 和 \(\xi_2\) 分别是 \(\mathbf{X}_1\) 和 \(\mathbf{X}_2\) 在 C 中的类，证明单生成理想 \(\mathrm{C}\xi_1\) 和 \(\mathrm{C}\xi_2\) 在 C 中的交不是有限生成的理想。*

\begin{enumerate}
\def\labelenumi{\arabic{enumi}.}
\setcounter{enumi}{8}
\item
  证明：在 A-模 E 中，若 \((\mathrm{F}_{n})_{n\geqslant1}\) 是有限生成子模的一个严格递增序列，则以诸 \(F_{n}\) 之并为 E 的子模的 F 不是有限生成的。由此推出，如果 E 是一个非有限生成的 A-模，则存在 E 的一个子模 F，使得 \(\gamma(\mathrm{F})=\aleph_{0}(=\mathrm{Card}(\mathbf{N}))\) .
\item
  设 E 是一个 Z-模，它是两个子模 M、N 的直和，且分别同构于 Z/2Z 和 Z/3Z；证明：不存在将 E 分解为两个非零子模的直和的其他分解方式。
\item
  设 A 是具有左分式域 (I, § 9, 习题 15) 的环。证明在 A-模 \(A_{s}\) 中，异于 \(A_{s}\) 和 \(\{0\}\) 的子模没有互补子模（参见 I，§ 9，习题 16）。
\item
  设 \(\mathbf{G}\) 为一个模，且 \(\mathbf{E}, \mathbf{F}\) 为两个子模，使得 \(\mathbf{E} \subset \mathbf{F}\) .
\end{enumerate}

\begin{enumerate}
\def\labelenumi{(\alph{enumi})}
\item
  若 F 是 G 的直因子，则 F/E 是 G/E 的直因子。若同时 E 是 F 的直因子，则 E 是 G 的直因子。
\item
  若E是G的一个直因子，则E也是F的一个直因子。若进一步F/E是G/E的一个直因子，则F是G的一个直因子。
\item
  给出两个子模 \(\mathbf{M}, \mathbf{N}\) （ \(\mathbf{Z}\) -模 \(\mathbf{E} = \mathbf{Z}^2\) 的）的例子，使得 \(\mathbf{M}\) 和 \(\mathbf{N}\) 都是 \(\mathbf{E}\) 的直因子，但 \(\mathbf{M} + \mathbf{N}\) 不是 \(\mathbf{E}\) 的直因子.
\item
  设 \(p\) 是一个素数， \(U_p\) 是 \(\mathbf{Z}\) -模 \(\mathbf{Q} / \mathbf{Z}\) 的子模，由模 \(\mathbf{Z}\) 的类（对应于形如 \(k / p^n\) 的有理数）组成 ( \(k \in \mathbf{Z}\) , \(n \in \mathbf{N}\) )。设 \(E\) 为积 \(\mathbf{Z}\) -模 \(M \times N\) ，其中 \(M\) 和 \(N\) 都同构于 \(U_p\) ，且 \(M\) 和 \(N\) 被典范地等同于 \(E\) 的子模。考虑自同态 \(u: (x, y) \mapsto (x, y + px)\) （ \(E\) 的）；证明 \(u\) 是双射。若 \(M' = u(M)\) ，证明 \(M\) 和 \(M'\) 都是 \(E\) 的直因子，但 \(M \cap M'\) 不是 \(E\) 的直因子（利用 (b)，并证明 \(Z\) -模 \(U_p\) 除自身和 \{0\} 外没有其他直因子）.
\end{enumerate}

¶ 13. 设 E 是一个 A-模，B 为环 End\_A E。证明对于元素 \(u \in B\) ，以下三个条件等价：（1）Bu 是左 B-模 B\_s 的一个直因子；（2）uB 是右 B-模 B\_d 的一个直因子；（3）Ker(u) 和 Im(u) 是 A-模 E 的直因子。（注意，作为 B\_s 的直因子的每个左理想都具有 Bp 的形式，其中 p 是 E 中的一个投影算子；参见 I，§ 8，习题 12，并利用 II，§ 1，第 9 号，命题 15 的推论 1 和 2。）

\begin{enumerate}
\def\labelenumi{\arabic{enumi}.}
\setcounter{enumi}{13}
\tightlist
\item
  设 L 是一个良序集，E 是一个 A-模，且 \((\mathbf{E}_{\lambda})_{\lambda \in L}\) 是 E 的子模的一个递增族，使得：（1）存在 \(\lambda \in L\) 使得 \(E_{\lambda} = \{0\}\) ； (2) E 是诸 \(E_{\lambda}\) （ \(\lambda \in L\) ）的并； (3) 如果 \(\lambda \in L\) 使得 \(\mu < \lambda\) 的集合具有最大元 \(\lambda'\) , \(E_{\lambda}\) 是 \(E_{\lambda'}\) 与一个子模 \(F_{\lambda}\) 的直和；（4）若 \(\lambda \in L\) 是 \(\mu < \lambda\) 的集合的上确界，则 \(E_{\lambda}\) 是诸 \(E_{\mu}\) （ \(\mu < \lambda\) ）的并。在这些条件下，证明 E 是诸 \(F_{\lambda}\) 的直和。（用超限归纳法证明每个 \(E_{\lambda}\) 是诸 \(F_{\mu}\) 的直和，其中 \(\mu \leqslant \lambda\) 。)
\end{enumerate}

¶ 15. 设 c 为无限基数。

\begin{enumerate}
\def\labelenumi{(\alph{enumi})}
\item
  设 I 为一个集合， \(R\{a, b\}\) 为 I 的元素之间的一种关系，使得对所有 \(\alpha \in I\) ，由 \(\beta \in I\) 中使 \(R\{\alpha, \beta\}\) 成立者组成的集合具有基数 \(\leqslant c\) 。证明存在一个良序集 L 和 I 的子集的一个递增族 \((\mathrm{I}_{\lambda})_{\lambda \in L}\) ，使得：(1) 存在 \(\lambda \in L\) 使得 \(I_{\lambda} = \varnothing\) ；(2) I 是诸 \(I_{\lambda}\) （ \(\lambda \in L\) ）的并； (3) 如果 \(\lambda \in L\) 使得 \(\mu < \lambda\) 的集合有最大元 \(\lambda'\) , \(\text{Card}(\mathbf{I}_{\lambda} - \mathbf{I}_{\lambda'}) \leqslant c\) ；（4）若 \(\lambda \in L\) 是 \(\mu < \lambda\) 的集合的上确界，则 \(I_{\lambda}\) 是诸 \(I_{\mu}\) （ \(\mu < \lambda\) ）的并；(5) 若 \(\alpha \in I_{\lambda}\) 且 \(R\{\alpha, \beta\}\) 成立，则 \(\beta \in I_{\lambda}\) 。（首先注意，对所有 \(\alpha \in I\) ，由满足如下条件的 \(\beta \in I\) 组成的集合——存在一个整数 n 和一个由 I 的元素组成的序列 \((\gamma_{j})_{1 \leqslant j \leqslant n}\) ，使得 \(\gamma_{1} = \alpha, \gamma_{n} = \beta\) 且 \(R\{\gamma_{j}, \gamma_{j+1}\}\) 对 \(1 \leqslant j \leqslant n - 1\) 成立——具有基数 \(\leqslant c\) 。然后通过超限归纳法构造诸 \(I_{\lambda}\) 。）
\item
  设 E 是一个 A-模，它是子模族 \((\mathbf{M}_{\alpha})_{\alpha\in I}\) 的直和，使得 \(\gamma(\mathbf{M}_{\alpha})\leqslant c\) 对所有 \(\alpha\in I\) 成立（参见习题 6）。设 f 是 E 的一个自同态。证明存在一个良序集 L 和 I 的子集的一个递增族 \((\mathbf{I}_{\lambda})_{\lambda\in L}\) ，满足 (a) 中的性质 (1)、(2)、(3) 和 (4)，并且此外，若记 \(\mathbf{E}_{\lambda} = \bigoplus_{\alpha \in \mathbf{I}_{\lambda}} \mathbf{M}_{\alpha}\) ，则 \(f(\mathbf{E}_{\lambda}) \subset \mathbf{E}_{\lambda}\) 。（将 (a) 应用于关系 \(R\{\alpha, \beta\}\) ：“存在 \(x \in \mathbf{M}_{\alpha}\) 使得 \(f(x)\) 在 \(\mathbf{M}_{\beta}\) 中的分量 \(\neq 0\)”。）对所有满足下列条件的 \(\lambda \in \mathbf{L}\) ：\(\mu < \lambda\) 的集合有最大元 \(\lambda'\) ，记 \(\mathbf{F}_{\lambda} = \bigoplus_{\alpha \in \mathbf{I}_{\lambda} - \mathbf{I}_{\lambda'}} \mathbf{M}_{\alpha}\) 。证明 \(\mathbf{E}\) 是诸 \(\mathbf{F}_{\lambda}\) 的直和（应用习题 14）。
\item
  在 (b) 的假设和记号下，进一步假设 \(f\) 是一个投影算子，并记 \(\mathbf{P} = f(\mathbf{E})\) 。设 \(\mathbf{P}_{\lambda} = \mathbf{P} \cap \mathbf{E}_{\lambda}\) ；证明如果 \(\lambda \in \mathbf{L}\) 使得 \(\mu < \lambda\) 的集合有最大元 \(\lambda'\) , \(\mathbf{P}_{\lambda'}\) 是 \(\mathbf{P}_{\lambda}\) 的一个直因子（参见习题 12）， \(\mathbf{P}_{\lambda}'\) 作为 \(\mathbf{P}_{\lambda'}\) 在 \(\mathbf{P}_{\lambda}\) 中的互补子模同构于 \(\mathbf{F}_{\lambda}\) 的一个直因子，且 \(\gamma(\mathbf{P}_{\lambda}') \leqslant c\) 。证明 \(\mathbf{P}\) 是诸 \(\mathbf{P}_{\lambda}'\) 的直和（应用习题 14）。
\item
  由 (c) 推出，若一个 A-模 E 是子模族 \((\mathrm{M}_{\alpha})_{\alpha\in I}\) 的直和，使得 \(\gamma(\mathrm{M}_{\alpha})\leqslant\mathfrak{c}\) ，那么 E 的每个直因子 P 也是子模族 \((\mathrm{N}_{\lambda})_{\lambda\in L}\) 的直和，使得 \(\gamma(\mathrm{N}_{\lambda})\leqslant\mathfrak{c}\) .
\end{enumerate}

\begin{enumerate}
\def\labelenumi{\arabic{enumi}.}
\setcounter{enumi}{15}
\tightlist
\item
  设 A 是一个环，使得存在一个 A-模 M，它有一个由 n 个元素组成的生成系，但包含一个由 \(n + 1\) 个元素组成的自由系。
\end{enumerate}

\begin{enumerate}
\def\labelenumi{(\alph{enumi})}
\item
  证明在 \(\mathbf{A}_s^n\) 中存在一个由 \(n + 1\) 个元素组成的自由系。
\item
  推断出在 M 中存在一个无限自由系（利用 (a) 通过归纳构造这样的系）。
\item
  设 C 为一个环，E 为自由 C-模 \(\mathrm{C}_{s}^{(\mathrm{N})}, (e_{n})\) 及其典范基，A 为 E 的自同态环。设 \(u_{1}\) 和 \(u_{2}\) 表示由下列条件定义的 E 的自同态： \(u_{1}(e_{2n}) = e_{n}, u_{1}(e_{2n+1}) = 0, u_{2}(e_{2n+1}) = e_{n}, u_{2}(e_{2n}) = 0\) 对所有 \(n \geqslant 0\) 。证明 \(u_{1}\) 和 \(u_{2}\) 构成 A-模 \(A_{s}\) 的一组基，并推断出在 A-模 \(A_{s}\) 中存在无限自由系。
\end{enumerate}

*(d) 设 A 是交换域上维数 \(\geqslant 2\) 的向量空间 E 的张量代数（参见 III，§ 5，第 1 号）。证明在 A-模 \(A_s\) 中存在无限自由系（注意，E 中两个线性无关的向量在 \(A_s\) 中也线性无关，并利用 \((b)\) ）。另一方面，对于每个整数 \(n > 0\) ，A-模 \(A_s^n\) 的每个基都有 \(n\) 个元素（参见 § 7，第 2 号，命题 3）。*

\begin{enumerate}
\def\labelenumi{\arabic{enumi}.}
\setcounter{enumi}{16}
\tightlist
\item
  设 A 是一个环。
\end{enumerate}

\begin{enumerate}
\def\labelenumi{(\alph{enumi})}
\item
  证明如果 A-模 \(\mathbf{A}_s^n\) 不同构于任何 \(\mathbf{A}_s^m\) （对于 \(m > n\) ），那么，对所有 \(p < n\) , \(\mathbf{A}_s^p\) 不同构于任何 \(\mathbf{A}_s^q\) （对于 \(q > p\) ）.
\item
  证明如果 A-模 \(A_{s}^{n}\) 不包含由 \(n + 1\) 个元素组成的自由系，则 A-模 \(A_{s}^{p}\) 不包含由 \(p + 1\) 个元素组成的自由系（对 p \textless{} n）.
\end{enumerate}

\begin{enumerate}
\def\labelenumi{\arabic{enumi}.}
\setcounter{enumi}{17}
\tightlist
\item
  设 A 是一个环，且存在一个整数 \(p\) 具有如下性质：对于每一个族 \((a_i)_{1 \leqslant i \leqslant p}\) （由 A 的 \(p\) 个元素组成），存在一个族 \((c_i)_{1 \leqslant i \leqslant p}\) （由 A 的元素组成），其中至少有一个元素不是右零因子，使得 \(\sum_{i=1}^{p} c_i a_i = 0\) .
\end{enumerate}

\begin{enumerate}
\def\labelenumi{(\alph{enumi})}
\item
  对 \(n\) 用归纳法证明：对于每个族 \((x_{i})\) （由 \(p^n\) 个 \(A_s^n\) 的元素组成），存在一个族 \((c_j)\) （由 \(p^n\) 个 \(A\) 的元素组成），其中至少有一个不是右零因子，使得 \(\sum_{j} c_j x_j = 0\) .
\item
  由 (a) 和习题 16 推出，对所有 n \textgreater{} 0，A-模 \(A_{s}^{n}\) 不包含由 \(n + 1\) 个元素组成的自由系。
\end{enumerate}

\begin{enumerate}
\def\labelenumi{\arabic{enumi}.}
\setcounter{enumi}{18}
\tightlist
\item
  设A是一个非零环且无零因子。
\end{enumerate}

\begin{enumerate}
\def\labelenumi{(\alph{enumi})}
\item
  证明对于 \(n > 1\) ，A-模 \(\mathbf{A}_s^n\) 不是单生成的。
\item
  给定一个整数 \(n \geqslant 1\) ，对于 A-模 \(\mathbf{A}_s^n\) 不包含由 \(n + 1\) 个元素组成的自由系，必要条件是 A 具有左分式域（I，§9，习题15）；反之，若此条件满足，则 \(\mathbf{A}_s^n\) 不包含由 \(n + 1\) 个元素组成的自由系，无论整数 \(n > 0\) 如何（利用习题 17 和 18 以及 I，§9，习题 16。）
\item
  证明：若 A 的每个左理想都是单生成的，则 A 具有左分式域（利用 (a) 以及 I，§9，习题 16）。
\end{enumerate}

\begin{enumerate}
\def\labelenumi{\arabic{enumi}.}
\setcounter{enumi}{19}
\tightlist
\item
  \begin{enumerate}
  \def\labelenumii{(\alph{enumii})}
  \tightlist
  \item
    设 A 是具有左分式域 (I, § 9, 习题 15) 的环。设 E 是一个 A-模；证明：若在 E 中 \((x_{i})_{1\leqslant i\leqslant r}\) 是一个自由系，且 \(y,z\) 为两个元素，使得 \(x_{1},\ldots ,x_{r},y\) 与 \(x_{1},\ldots ,x_{r},z\) 都是相关系，那么 \(x_{2},\ldots ,x_{r},y,z\) 是一个相关系。
  \end{enumerate}
\end{enumerate}

\begin{enumerate}
\def\labelenumi{(\alph{enumi})}
\setcounter{enumi}{1}
\tightlist
\item
  设 A 为商环 Z/6Z，并设 \((e_{1}, e_{2})\) 为 \(A^{2}\) 的典范基；证明如果 \(a = 2e_{1} + 3e_{2}\) ，则 a 与 \(e_{1}\) 构成一个相关系，a 与 \(e_{2}\) 亦然，尽管 \(e_{1}\) 和 \(e_{2}\) 构成一个自由系。
\end{enumerate}

\begin{enumerate}
\def\labelenumi{\arabic{enumi}.}
\setcounter{enumi}{20}
\item
  设 A 为一个环，b、c 为 A 中两个非右零因子的元素，M 为 A-模 A/Abc，N 为 M 的子模 Ac/Abc。要使 N 成为 M 的一个直因子，当且仅当存在 A 中的两个元素 x、y，使得 \(xb + cy = 1\) 。（若 \(\varepsilon\) 是 1 在 M 中的类，证明 N 在 M 中的一个互补子空间由形如 \((1 - yc)\varepsilon\) 的元素生成，其零化子为理想 Ac。）
\item
  若一个A-模E存在一个以集合I为指标集的基，且A与I两个集合中至少有一个是无限的，证明E与 \(A \times I\) 等势.
\item
  \begin{enumerate}
  \def\labelenumii{(\alph{enumii})}
  \tightlist
  \item
    设 M、N 是 A 模 E 的两个子集，m 和 n 分别是它们的零化子；证明 M ∩ N 的零化子包含 m + n，并给出一个例子说明它不同于 m + n。
  \end{enumerate}
\end{enumerate}

\begin{enumerate}
\def\labelenumi{(\alph{enumi})}
\setcounter{enumi}{1}
\item
  在积模 \(\prod_{i} \mathbf{E}_{i}\) 中，子集 \(\mathbf{F}\) 的零化子是它的各投影的零化子的交。
\item
  在自由 A-模中，元素 \(\neq0\) 的零化子仅包含 A 中 0 的左因子；特别地，若 A 是无零因子环，则自由模的每个元素 \(\neq0\) 都是自由的。
\end{enumerate}

\begin{enumerate}
\def\labelenumi{\arabic{enumi}.}
\setcounter{enumi}{23}
\tightlist
\item
  设 E 是一个 A-模，M 和 N 是 E 的两个子模。M 到 N 的传送子，记为 N:M，是由 \(a \in A\) 中使 \(aM \subset N\) 成立者组成的集合；它是 A 的一个双边理想，等于 \(\text{Ann}((\mathbf{M} + \mathbf{N})/\mathbf{N})\) .
\end{enumerate}

\begin{enumerate}
\def\labelenumi{(\alph{enumi})}
\item
  设 A 是一个交换环，M 是一个 A-模，N 是 M 的一个子模，x 是 M 的一个元素；证明 \(N \cap Ax = (N:Ax)x\) .
\item
  设 A 是一个交换环，且 \(a, b\) 为 A 的两个理想。证明 \(a:b \supset a\) ，并且 \((a:b)/a\) 是同构于 \(\operatorname{Hom}(A/b, A/a)\) 的一个 A-模.
\end{enumerate}

\begin{enumerate}
\def\labelenumi{\arabic{enumi}.}
\setcounter{enumi}{24}
\item
  设 E, F 为两个 A-模，且 \(u: E \to F\) 为一个线性映射。证明映射 \((x, y) \mapsto (x, y - u(x))\) （积模 \(E \times F\) 到自身的）是 \(E \times F\) 的一个自同构。由此推出，如果存在一个线性映射 \(v: F \to E\) 和一个 \(a \in E\) ，使得 \(v(u(a)) = a\) ，则存在 \(E \times F\) 的一个自同构 w，使得 \(w(a, 0) = (0, u(a))\) .
\item
  设 E 是一个自由 A-模，其基至少包含两个元素。
\end{enumerate}

\begin{enumerate}
\def\labelenumi{(\alph{enumi})}
\item
  证明：E到自身的每一个半线性映射（相对于A的一个自同构），若它与E的所有自同构可交换，则它必为一个位似映射 \(x \mapsto \alpha x (\alpha \in \mathbf{A})\) .
\item
  由 (a) 推出， \(\operatorname{End}_{\mathbf{A}}(\mathbf{E})\) 的中心是中心位似环，同构于 A 的中心。
\item
  由 (a) 推出，群 \(\mathbf{GL}(\mathbf{E})\) 的中心是 E 的可逆中心位似的群，它同构于 A 的中心中可逆元素组成的乘法群。
\end{enumerate}

\begin{enumerate}
\def\labelenumi{\arabic{enumi}.}
\setcounter{enumi}{26}
\tightlist
\item
  设 A 是一个交换环。证明若 A 的一个理想 \(\mathfrak{J}\) 是自由 A-模，则 \(\mathfrak{J}\) 是单生成 A-模（换言之，是一个主理想）。给出一个例子，说明该命题不能推广到非交换环的左理想（习题 16）。
\end{enumerate}

